\documentclass{dlproblemset}

\title{LSTMs e GRUs}
\subtitle{Lista de Exercícios}

\begin{document}

\subsection*{Questões de Múltipla Escolha}

\begin{enumerate}[I)]

% ------------------------------------------------------------------
% Q1 — Medium (contagem de parâmetros LSTM; distrator = GRU / sem bias)
% ------------------------------------------------------------------
\item Uma célula LSTM recebe entradas de dimensão $is = 10$ e mantém um estado oculto de
dimensão $hs = 20$. Cada uma das quatro camadas densas internas tem pesos
$\mathbf{\Theta}_\bullet \in \mathbb{R}^{(hs + is) \times hs}$ e \textit{bias}
$\mathbf{b}_\bullet \in \mathbb{R}^{hs}$. Quantos parâmetros treináveis (pesos e \textit{bias})
a célula possui no total?

\begin{enumerate}[a)]
    \item $2480$.
    \item $2400$.
    \item $1860$.
    \item $1240$.
\end{enumerate}
% R: a

% ------------------------------------------------------------------
% Q2 — Medium (comportamento das gates na atualização de c_t)
% ------------------------------------------------------------------
\item Considere a atualização da memória $\mathbf{c}_t = \mathbf{f}_t \odot \mathbf{c}_{t-1} + \mathbf{i}_t \odot \tilde{\mathbf{c}}_t$.
Qual configuração das \textit{gates} faz a célula \textbf{preservar integralmente} a memória
de longo prazo e \textbf{ignorar} a entrada do passo atual?

\begin{enumerate}[a)]
    \item $\mathbf{f}_t \to 1$ e $\mathbf{i}_t \to 0$.
    \item $\mathbf{f}_t \to 0$ e $\mathbf{i}_t \to 1$.
    \item $\mathbf{f}_t \to 1$ e $\mathbf{i}_t \to 1$.
    \item $\mathbf{f}_t \to 0$ e $\mathbf{i}_t \to 0$.
\end{enumerate}
% R: a

% ------------------------------------------------------------------
% Q3 — Medium (cálculo de h_t; distratores = erros comuns)
% ------------------------------------------------------------------
\item Em um passo de uma LSTM o \textit{output gate} resultou em $o_t = 0.6$ e a memória
atualizada foi $c_t = 0.5$. Usando $h_t = o_t \odot \tanh(c_t)$, qual é o valor de $h_t$?

\begin{enumerate}[a)]
    \item $\approx 0.2773$.
    \item $\approx 0.3000$.
    \item $\approx 0.2913$.
    \item $\approx 0.4621$.
\end{enumerate}
% R: a

% ------------------------------------------------------------------
% Q4 — Medium (update gate da GRU; direção de z_t)
% ------------------------------------------------------------------
\item Na GRU, o estado oculto é atualizado por
$\mathbf{h}_t = (1 - \mathbf{z}_t) \odot \mathbf{h}_{t-1} + \mathbf{z}_t \odot \tilde{\mathbf{h}}_t$. Quando o \textit{update gate}
satisfaz $\mathbf{z}_t \to 0$, o que ocorre com o estado oculto?

\begin{enumerate}[a)]
    \item $\mathbf{h}_t \approx \mathbf{h}_{t-1}$.
    \item $\mathbf{h}_t \approx \tilde{\mathbf{h}}_t$.
    \item $\mathbf{h}_t \approx 0$.
    \item $\mathbf{h}_t \approx \mathbf{r}_t \odot \mathbf{h}_{t-1}$.
\end{enumerate}
% R: a

\end{enumerate}

\subsection*{Gabarito - Objetivas}
\color{blue}
\begin{enumerate}[I)]
    \item \textbf{a.} Cada \textit{gate} tem $(hs+is)\cdot hs$ pesos mais $hs$ \textit{bias},
    isto é, $(30)(20) + 20 = 620$ parâmetros. Com quatro camadas densas,
    $4 \cdot 620 = 2480$. A alternativa (b) $2400$ esquece os \textit{bias}
    ($4 \cdot 600$); (c) $1860$ usa apenas três camadas (é a conta de uma GRU); (d) $1240$
    conta apenas duas camadas.
    \item \textbf{a.} Para manter $\mathbf{c}_{t-1}$ e bloquear a contribuição nova é preciso
    $\mathbf{f}_t \to 1$ (preserva a memória) e $\mathbf{i}_t \to 0$ (zera $\mathbf{i}_t \odot \tilde{\mathbf{c}}_t$). As demais
    ou apagam a memória ($\mathbf{f}_t \to 0$) ou deixam a entrada atualizar o estado ($\mathbf{i}_t \to 1$).
    \item \textbf{a.} $h_t = 0.6 \cdot \tanh(0.5) = 0.6 \cdot 0.4621 \approx 0.2773$.
    A alternativa (b) $0.3000$ esquece o $\tanh$; (c) $0.2913$ aplica $\tanh$ \emph{depois}
    do produto, $\tanh(0.3)$; (d) $0.4621$ esquece de multiplicar por $o_t$.
    \item \textbf{a.} Com $\mathbf{z}_t \to 0$, $\mathbf{h}_t = 1 \cdot \mathbf{h}_{t-1} + 0 \cdot \tilde{\mathbf{h}}_t = \mathbf{h}_{t-1}$:
    o estado anterior é copiado. A alternativa (b) corresponde a $\mathbf{z}_t \to 1$; (c) e (d) não
    correspondem à equação de atualização.
\end{enumerate}
\color{black}

\newpage

\subsection*{Questões Dissertativas}

\begin{enumerate}[I)]

% ------------------------------------------------------------------
% D1 — Medium (forward pass escalar passo a passo)
% ------------------------------------------------------------------
\item Considere uma célula LSTM \emph{escalar} com estado inicial $c_0 = 0.2$,
$h_0 = 0.5$ e entrada $x_0 = -0.4$, de modo que o vetor concatenado é
$[h_0,\, x_0] = [0.5,\, -0.4]$. Cada camada densa interna tem pesos
$\theta_\bullet = [\theta_h,\, \theta_x]$ que multiplicam $[h_{t-1},\, x_t]^T$:
\[
    \theta_f = [0.4,\, 0.6],\ b_f = -0.1; \qquad
    \theta_i = [0.5,\, -0.2],\ b_i = 0.1;
\]
\[
    \theta_c = [-0.3,\, 0.5],\ b_c = 0.4; \qquad
    \theta_o = [0.2,\, 0.3],\ b_o = 0.
\]
Calcule, passo a passo (use 4 casas decimais):
\begin{enumerate}[(a)]
    \item o \textit{forget gate} $f_t$ e o \textit{input gate} $i_t$;
    \item o candidato $\tilde{c}_t$;
    \item a nova memória $c_t = f_t\, c_0 + i_t\, \tilde{c}_t$;
    \item o \textit{output gate} $o_t$ e o estado oculto $h_t = o_t\, \tanh(c_t)$.
\end{enumerate}

% ------------------------------------------------------------------
% D2 — Hard (porquê do gradiente; paralelo LSTM/GRU)
% ------------------------------------------------------------------
\item A LSTM foi projetada para mitigar o \textit{vanishing gradient} das RNNs simples.
\begin{enumerate}[(a)]
    \item A partir da atualização $\mathbf{c}_t = \mathbf{f}_t \odot \mathbf{c}_{t-1} + \mathbf{i}_t \odot \tilde{\mathbf{c}}_t$, mostre que,
    pelo \emph{caminho direto} no tempo, $\dfrac{\partial \mathbf{c}_t}{\partial \mathbf{c}_{t-1}} = \mathbf{f}_t$.
    Explique por que isso evita o desaparecimento do gradiente, comparando com a RNN
    \textit{vanilla}, cujo gradiente acumula o produto
    $\prod_k \mathbf{\Theta}_{hh}\,(1 - \tanh^2(\mathbf{z}_k))$.
    \item \emoji{skull-and-crossbones} A GRU não mantém um estado de memória $\mathbf{c}_t$ separado;
    o estado é atualizado por $\mathbf{h}_t = (1 - \mathbf{z}_t)\, \mathbf{h}_{t-1} + \mathbf{z}_t\, \tilde{\mathbf{h}}_t$. Identifique qual
    termo desempenha, na GRU, o papel análogo ao do \textit{forget gate} da LSTM na preservação
    do gradiente e diga sob qual condição do \textit{update gate} $\mathbf{z}_t$ o gradiente é carregado
    quase sem atenuação.
\end{enumerate}

\end{enumerate}

\subsection*{Gabarito}
\color{blue}
\begin{enumerate}[I)]

\item \textit{Forward pass} escalar de um passo da LSTM.
\begin{enumerate}[(a)]
    \item \textit{Forget} e \textit{input gates}:
    \begin{align*}
        f_t &= \sigma\bigl(0.4 \cdot 0.5 + 0.6 \cdot (-0.4) - 0.1\bigr)
            = \sigma(-0.14) \approx 0.4651 \\
        i_t &= \sigma\bigl(0.5 \cdot 0.5 + (-0.2) \cdot (-0.4) + 0.1\bigr)
            = \sigma(0.43) \approx 0.6059
    \end{align*}
    \item Candidato:
    \[
        \tilde{c}_t = \tanh\bigl(-0.3 \cdot 0.5 + 0.5 \cdot (-0.4) + 0.4\bigr)
            = \tanh(0.05) \approx 0.0500
    \]
    \item Nova memória:
    \[
        c_t = 0.4651 \cdot 0.2 + 0.6059 \cdot 0.0500
            \approx 0.0930 + 0.0303 = 0.1233
    \]
    \item \textit{Output gate} e estado oculto:
    \begin{align*}
        o_t &= \sigma\bigl(0.2 \cdot 0.5 + 0.3 \cdot (-0.4) + 0\bigr)
            = \sigma(-0.02) \approx 0.4950 \\
        h_t &= 0.4950 \cdot \tanh(0.1233) \approx 0.4950 \cdot 0.1227 \approx 0.0607
    \end{align*}
\end{enumerate}

\item Por que o caminho da memória preserva o gradiente.
\begin{enumerate}[(a)]
    \item Derivando $\mathbf{c}_t = \mathbf{f}_t \odot \mathbf{c}_{t-1} + \mathbf{i}_t \odot \tilde{\mathbf{c}}_t$ em relação a $\mathbf{c}_{t-1}$,
    o termo direto $\mathbf{f}_t \odot \mathbf{c}_{t-1}$ contribui com $\mathbf{f}_t$; os fatores $\mathbf{f}_t$, $\mathbf{i}_t$ e
    $\tilde{\mathbf{c}}_t$ dependem de $\mathbf{c}_{t-1}$ apenas indiretamente, através de $\mathbf{h}_{t-1}$, formando
    caminhos secundários menores. Logo, pelo caminho direto,
    $\partial \mathbf{c}_t / \partial \mathbf{c}_{t-1} = \mathbf{f}_t$. Ao longo de $T$ passos esse caminho acumula
    $\prod_k \mathbf{f}_k$: são multiplicações por portas aprendidas e \textbf{somas}, sem $\tanh'$
    repetido. Quando o \textit{forget gate} aprende $\mathbf{f}_k \approx 1$, o produto permanece
    $\approx 1$ e o gradiente é preservado (\textit{constant error carousel}). Na RNN
    \textit{vanilla}, cada passo multiplica por $\mathbf{\Theta}_{hh}\,(1 - \tanh^2(\mathbf{z}_k))$; como
    $1 - \tanh^2 \le 1$, o produto tende a $0$ (\textit{vanishing}) ou, com $|\mathbf{\Theta}_{hh}|$
    grande, explode.
    \item Na GRU, o fator $(1 - \mathbf{z}_t)$ multiplica $\mathbf{h}_{t-1}$ no caminho direto, então
    $\partial \mathbf{h}_t / \partial \mathbf{h}_{t-1} = (1 - \mathbf{z}_t)$ por esse caminho, papel análogo ao $\mathbf{f}_t$ da
    LSTM. Quando $\mathbf{z}_t \approx 0$ (\textit{update gate} ``fechado''), tem-se $(1 - \mathbf{z}_t) \approx 1$
    e o estado é carregado com gradiente $\approx 1$, preservando a informação ao longo do tempo.
\end{enumerate}

\end{enumerate}

\end{document}
